\Needspace{5\baselineskip}\pdfbookmark[2]{Second division}{witwen-II-2}
\begin{center}\textit{Second division.}\end{center}
\noindent\hspace*{\parindent}
Since I have already described in detail, in my opinion of 8 January of this year, the paths to be followed in drawing up the balance sheet of the widows' fund, and found no reason for alterations when carrying out the work, I need only refer to that opinion. I shall merely remark here that the essence of the method consists in determining the present monetary value of the fund's obligations under the three headings
\begin{quote}
Obligations towards the present 19 widows,\par
Obligations towards the widows (and orphans) of the present 51 members,\par
Obligations towards the widows and orphans of those joining in the future,
\end{quote}
\noindent where, under the second and third headings, the present monetary value of the contributions made in return is to be deducted. I shall first give some general explanations here.

The date to which all future payments are referred by discounting has been chosen as 1 October 1845. It is assumed that the widows and members have already respectively received their pensions and paid their contributions falling due by that date.
\par\SourceLeaf{156}
For the calculation I have used the highly valuable mortality tables that \textsc{Brune} derived from the experience of 31500 married couples in the Prussian widows' fund. These tables are deficient only for the deaths of men at the highest ages, since the fund's records contained insufficient data for this purpose. I have therefore preferred to reckon deaths among men over 80 according to the same proportions that the tables give for the female sex. To calculate the modifications in the value of widows' pensions arising from taking minor children into account, where such children are present, I have used \textsc{Deparcieux}'s mortality tables for the children. The choice of tables, both for the mortality of children and for that of men at the highest ages, has, however, only a very small influence on the results of the calculation.

All the results depend essentially on the rate of interest on which they are based. To leave nothing to be desired here, I have therefore carried out all the calculations both at a rate of 4 per cent and at one of \(3\tfrac{1}{2}\) per cent, although this exactly doubled the work. The currency amounts are always to be understood as gold. Throughout, the calculation has been carried to fractions of a thaler, but these fractions have been omitted in the present copy. This may here and there produce a difference of one or a couple of units in the sums. I mention this minor circumstance here only so that anyone who might check one or other sum should not take such differences as signs of inaccurate calculation, since, quite the reverse, they result from the greater precision observed in the calculation.

\Needspace{6\baselineskip}\pdfbookmark[3]{I. Calculation of the monetary value of widows' pensions}{witwen-II-2-I}
\par\medskip\noindent I.\quad Calculation of the monetary value, discounted to 1 October 1845, of widows' pensions.\par\smallskip
The total number of cases in which pensions have been granted since the establishment of the widows' fund until now is 68. To make my work as easy to follow as possible, I have given them consecutive numbers. The figures in the following table express the sums that would have to be paid to the individual widows if their claims were to be settled, or, conversely, the sums the widows would have to pay in if they wished to purchase an entitlement such as they now possess. The calculation is based on the present pension rate of 250 Rthl., except for No. 58, where it amounts to 200 Rthl. The widows 50. 53. 58. 64. 68 have children under 20 years of age, who have been taken exactly into account in the calculation, as has the fact that pensions are always still paid for the full month in which a recipient leaves.

\begin{center}
\footnotesize\setlength{\tabcolsep}{3pt}
\begin{tabular}{r|c|r|r||r|c|r|r}
\multicolumn{2}{c|}{}&\multicolumn{2}{c||}{\shortstack{Present value\\of the pension.}}&\multicolumn{2}{c|}{}&\multicolumn{2}{c}{\shortstack{Present value\\of the pensions.}}\\\cline{3-4}\cline{7-8}
No.& &\(3\tfrac{1}{2}\) per cent&4 per cent&No.& &\(3\tfrac{1}{2}\) per cent&4 per cent\\\hline
29&L.&957 Thl.&941 Thl.&58&G.&3172 Thl.&3000 Thl.\\
31&S.&1756 ,, &1708 ,, &59&W.&2850 ,, &2732 ,, \\
41&W.&1122 ,, &1101 ,, &60&S.&2606 ,, &2507 ,, \\
44&F.&2543 ,, &2448 ,, &62&B.&1648 ,, &1605 ,, \\
50&H.&3987 ,, &3759 ,, &63&G.&2113 ,, &2045 ,, \\
51&S.&735 ,, &725 ,, &64&M.&4170 ,, &3932 ,, \\
52&P.&2352 ,, &2270 ,, &65&H.&3026 ,, &2894 ,, \\
53&M.&2789 ,, &2676 ,, &66&H.&1916 ,, &1859 ,, \\
55&H.&2813 ,, &2698 ,, &68&M.&4175 ,, &3936 ,, \\\cline{5-8}
57&S.&1798 ,, &1748 ,, &\multicolumn{2}{c}{Total}&46529 Thl.&44582 Thl.
\end{tabular}
\end{center}
